\documentclass[hu]{../../elegantbook}
% ========== Borító és alapinformációk ==========
\title{ElegantBook teszdokumentum}
\subtitle{A sablon alapfunkciói}
\author{Tesztmérnök}
\institute{Sablon Tesztcsoport}
\date{\today}
\version{v1.0}
\extrainfo{A dokumentum kizárólag sablonfunkció tesztelésére szolgál, nincs tudományos értéke}
\bioinfo{Tesztköteg}{TEST-FULL-20260506}
\logo{../../figure/logo-blue.png}
\cover{../../figure/cover.jpg}

% ========== Tartalomjegyzék és alapbeállítások ==========
\setcounter{tocdepth}{3} % Tartalomjegyzék megjelenítése 3 szintig

% ========== Teszt irodalomjegyzék ==========
\addbibresource[location=local]{../en/reference.bib}

\begin{document}

\maketitle
\frontmatter
\tableofcontents

% ========== Előszó (szakaszokra bontva) ==========
\pagenumbering{Roman}
\setcounter{page}{0}
\chapter*{Tesztelőszó}
\markboth{Tesztelőszó}{}
\section*{Dokumentum leírása}
\markright{Dokumentum leírása}
Ez a dokumentum az ElegantBook teljes funkcióinak tesztesete, lefedi a sablon összes alap és haladó funkcióját. Minden tartalom kitalált tesztadat, gyakorlati jelentősége nincs.
\newpage
\section*{Teszt hatókör}
\markright{Teszt hatókör}
A teszt tartalmazza: borító, tartalomjegyzék, betűtípusok, többszintű listák, összetett matematikai képletek, összes matematikai környezet, kereszthivatkozások, ábrák és táblázatok, margójegyzetek, fejezetösszegzések, feladatok, irodalomjegyzék, függelékek.

\mainmatter

% ========== 1. fejezet: Alapformázás és többszintű listák ==========
\chapter{Alapformázás és lista teszt}
\begin{introduction}[A fejezet tesztpontjai]
\item Kínai betűtípusok és szövegformázás
\item 3 szintű rendezetlen lista beágyazása
\item 3 szintű rendezett lista beágyazása
\item Szokásos bekezdések és hangsúlyozott szöveg
\end{introduction}

\section{Betűtípus és szöveg teszt}
Szokásos főszöveg teszt. \textbf{Félkövér}, \textit{Dőlt}, \underline{Aláhúzott}, \textcolor{red}{Piros szöveg teszt}.

\section{3 szintű rendezetlen lista teszt}
\begin{itemize}
\item 1. szintű teszt elem A
\item 1. szintű teszt elem B
    \begin{itemize}
    \item 2. szintű teszt elem B1
    \item 2. szintű teszt elem B2
        \begin{itemize}
        \item 3. szintű teszt elem B2-1
        \item 3. szintű teszt elem B2-2
            \begin{itemize}
            \item 4. szintű teszt elem B2-1a
            \item 4. szintű teszt elem B2-2b
            \end{itemize}
        \end{itemize}
    \end{itemize}
\end{itemize}

\section{3 szintű rendezett lista teszt}
\begin{enumerate}
\item 1. szintű rendezett teszt 1
\item 1. szintű rendezett teszt 2
    \begin{enumerate}
    \item 2. szintű rendezett teszt 2-1
    \item 2. szintű rendezett teszt 2-2
        \begin{enumerate}
        \item 3. szintű rendezett teszt 2-2-1
        \item 3. szintű rendezett teszt 2-2-2
            \begin{enumerate}
            \item 4. szintű rendezett teszt 2-2-2-1
            \item 4. szintű rendezett teszt 2-2-2-2
            \end{enumerate}
        \end{enumerate}
    \end{enumerate}
\end{enumerate}

% ========== 2. fejezet: Összetett matematikai képletek tesztje ==========
\chapter{Összetett matematikai képletek tesztje}
\begin{introduction}
\item Egysoros összetett képlet
\item Többsoros képletek (align)
\item Mátrixok / Determinánsok
\item Többszörös integrálok / Összegzések / Határértékek / Törtek
\item Képletcímkék és kereszthivatkozások
\end{introduction}

\section{Egysoros összetett képlet}
Sorba illesztett összetett képletek: $\sum_{i=1}^n \frac{x_i^2}{\sqrt{y_i+1}} \geq \bar{x}^2$, $\lim_{n\to\infty} \left(1+\frac{1}{n}\right)^n = e$.

\section{Többsoros képletek és mátrixok}
\begin{align}
a^2 + b^2 &= c^2 \label{eq:gougu} \\
\int_{0}^{1}\int_{0}^{1} f(x,y) dxdy &= 1 \label{eq:double} \\
\begin{vmatrix}
a & b \\
c & d
\end{vmatrix} &= ad-bc \label{eq:det}
\end{align}

\section{Nagy operátoros képletek}
\begin{equation}\label{eq:series}
\sum_{k=0}^{\infty} \frac{x^k}{k!} = e^x,\quad \oint_{C} Pdx+Qdy = \iint_{D}\left(\frac{\partial Q}{\partial x}-\frac{\partial P}{\partial y}\right)dxdy
\end{equation}

Képlet hivatkozás: Pitagorasz-tétel \eqref{eq:gougu}, kettős integrál \eqref{eq:double}, determináns \eqref{eq:det}, sorozat \eqref{eq:series}.

% ========== 3. fejezet: Összes matematikai környezet tesztje ==========
\chapter{Összes matematikai környezet teljes tesztje}
\begin{introduction}
\item Definíció / Tétel / Lemma / Következtetés / Axióma / Posztulátum
\item Állítás / Példa / Probléma / Feladat
\item Megjegyzés / Észrevétel / Következtetés / Feltételezés / Tulajdonság
\item Megoldás / Bizonyítási környezet
\end{introduction}

\section{Környezet teszt}

% 1. Definíció
\ifdefined\simpleTest
\begin{definition}[Teszt definíció]\label{def:test}
\else
\begin{definition}[Teszt definíció]{test}
\fi
Ez a definíciós környezet teszt tartalma, nincs valós matematikai jelentése, csak stílus ellenőrzésére szolgál.
\end{definition}

% 2. Tétel
\ifdefined\simpleTest
\begin{theorem}[Teszt tétel]\label{thm:test}
\else
\begin{theorem}[Teszt tétel]{test}
\fi
Ha $x>0$, akkor $x+1>1$, megfelel a \eqref{eq:gougu} képletnek.
\end{theorem}

% 3. Lemma
\ifdefined\simpleTest
\begin{lemma}[Teszt lemma]\label{lem:test}
\else
\begin{lemma}[Teszt lemma]{test}
\fi
Bármely valós $x$ esetén teljesül $x^2\geq0$.
\end{lemma}

% 4. Következtetés
\ifdefined\simpleTest
\begin{corollary}[Teszt következtetés]\label{cor:test}
\else
\begin{corollary}[Teszt következtetés]{test}
\fi
A \ref{lem:test} lemmából következik: $(x-1)^2\geq0$.
\end{corollary}

% 5. Axióma
\ifdefined\simpleTest
\begin{axiom}[Teszt axióma]\label{axi:test}
\else
\begin{axiom}[Teszt axióma]{test}
\fi
Teszt axióma, gyakorlati jelentősége nincs.
\end{axiom}

% 6. Posztulátum
\ifdefined\simpleTest
\begin{postulate}[Teszt posztulátum]\label{pos:test}
\else
\begin{postulate}[Teszt posztulátum]{test}
\fi
Teszt posztulátum, csak környezet ellenőrzésére.
\end{postulate}

% 7. Állítás
\ifdefined\simpleTest
\begin{proposition}[Teszt állítás]\label{pro:test}
\else
\begin{proposition}[Teszt állítás]{test}
\fi
Teszt állítás tartalom, kapcsolódik a \ref{def:test} definícióhoz.
\end{proposition}

% 8. Példa
\begin{example}\label{exa:test}
Példa környezet teszt, automatikus számozás.
\end{example}

% 9. Probléma
\begin{problem}\label{probl:test}
Probléma környezet teszt, megoldandó kérdések felvetésére.
\end{problem}

% 10. Feladat
\begin{exercise}\label{exer:test}
Számítsd ki: $1+2+3+\dots+10$.
\end{exercise}

% 11. Megjegyzés
\begin{note}
Megjegyzés környezet: nincs számozás, fontos emlékeztetőkhöz.
\end{note}

% 12. Észrevétel
\begin{remark}
Észrevétel környezet: kiegészítő magyarázatokhoz.
\end{remark}

% 13. Következtetés
\begin{conclusion}
Következtetés környezet: összefoglaló teszt tartalom.
\end{conclusion}

% 14. Feltételezés
\begin{assumption}
Feltételezés környezet: teszt feltételek beállítása.
\end{assumption}

% 15. Tulajdonság
\begin{property}
Tulajdonság környezet: teszt objektum jellemzői.
\end{property}

% 16. Megoldás
\begin{solution}
A \ref{exer:test} feladat megoldása: Az összeg $55$ (teszt válasz).
\end{solution}

% 17. Bizonyítás
\begin{proof}
A \ref{thm:test} tétel bizonyítása: Mivel $x>0$, mindkét oldalhoz 1-et adva a tétel bizonyított.
\end{proof}

% ========== 4. fejezet: Kereszthivatkozások · Ábrák · Margójegyzetek ==========
\chapter{Kereszthivatkozások és bővített funkciók}
\begin{introduction}
\item Minden típusú kereszthivatkozás
\item Egyszerű ábra és táblázat teszt
\item Margójegyzet funkció teszt
\item Fejezet navigáció ellenőrzése
\end{introduction}

\section{Összes kereszthivatkozás összefoglalása}
Definíció \ref{def:test}, Tétel \ref{thm:test}, Lemma \ref{lem:test}, Következtetés \ref{cor:test}, Axióma \ref{axi:test}, Posztulátum \ref{pos:test}, Állítás \ref{pro:test}, Példa \ref{exa:test}, Probléma \ref{probl:test}, Feladat \ref{exer:test}.

\section{Ábra és táblázat teszt}
\begin{figure}[h]
    \centering
    \includegraphics[width=0.5\textwidth]{../../image/scatter.jpg}
    \caption{Teszt ábra}\label{fig:test}
\end{figure}
\begin{table}[h]
    \centering
    \begin{tabular}{c*{6}{c}}
    \toprule
        $n$ & 1 & 2 & 3 & 4 & 5 & 6 \\
    \midrule
        $a_n$ & 1 & 1 & 2 & 3 & 5 & 8 \\ 
    \bottomrule
    \end{tabular}
    \caption{Teszt táblázat}\label{tab:test}
\end{table}
Ábra/táblázat hivatkozás: \figref{fig:test} a teszt ábra, \tabref{tab:test} a teszt táblázat.

% ========== Fejezetvégi feladatok ==========
\begin{problemset}[Fejezet teszt feladatok]
\item Ellenőrizd a \ref{def:test} definíció megjelenítési stílusát
\item Vedd le a \eqref{eq:double} képletet
\item Ellenőrizd, hogy a \ref{fig:test} ábra helyesen jelenik-e meg
\end{problemset}

% ========== Irodalomjegyzék ==========
\nocite{*}
\printbibliography

% ========== Függelék ==========
\appendix
\chapter{Teszt függelék}
A függelék tartalma a függelék fejezetek, fejlécék, oldalszámozás és formázási stílus ellenőrzésére szolgál, nincs valós információja.
\section{Teszt függelék 1}
A függelék tartalma a függelék fejezetek, fejlécék, oldalszámozás és formázási stílus ellenőrzésére szolgál, nincs valós információja.
\newpage
\section{Teszt függelék 2}
A függelék tartalma a függelék fejezetek, fejlécék, oldalszámozás és formázási stílus ellenőrzésére szolgál, nincs valós információja.
\end{document}